\chapter{Pulmonary Function Tests}

\section*{What Is This Test?}
Pulmonary function tests measure airflow, lung volumes, and gas transfer. Spirometry is the core test and records forced vital capacity and forced expiratory volume in one second. Additional measurements may include lung volumes and diffusing capacity for carbon monoxide.

\section*{Alternative Names}
PFTs; Lung Function Tests; Spirometry; Pulmonary Function Study; Respiratory Function Testing.

\section*{Why This Test Is Ordered}
PFTs help diagnose and classify asthma, chronic obstructive pulmonary disease, restrictive lung disease, unexplained breathlessness, occupational lung disease, and impairment before or during selected treatments. They can assess response to bronchodilator medication and track change over time.

\section*{How This Test Is Performed}
The patient breathes through a mouthpiece while wearing a nose clip and follows coached instructions for slow and forceful breaths. Spirometry is repeated for quality and reproducibility. A bronchodilator may be given and the test repeated. Lung volumes and gas transfer require additional maneuvers on specialized equipment.

\section*{Preparation, Risks, and Limitations}
Follow instructions about smoking, heavy meals, vigorous exercise, and inhalers. Wear loose clothing and tell the laboratory about recent surgery, chest pain, infection, or inability to perform forceful breathing. Temporary cough, dizziness, or breathlessness may occur. Poor effort and suboptimal technique can make results uninterpretable.

\section*{Understanding Results}
An obstructive pattern is characterized by reduced expiratory flow relative to vital capacity; restriction is suggested by reduced total lung capacity and requires confirmation with lung-volume testing. Bronchodilator response, diffusion capacity, and clinical context help distinguish causes. A normal test does not exclude intermittent asthma or all forms of lung disease.

\section*{References}
American Thoracic Society and European Respiratory Society standards for pulmonary-function testing.

\clearpage
\section*{Understanding Results and Follow-up}
The laboratory report should identify the specimen, method, units, reference interval or decision limit, and whether the result is positive, negative, abnormal, or inconclusive. Reference intervals vary with age, sex, pregnancy, specimen type, assay, and laboratory; a result outside a range is not automatically a diagnosis, and a result within a range does not always exclude disease. Discuss unexpected results with the ordering clinician, who can decide whether repeat or confirmatory testing, treatment, or referral is appropriate. Do not change medicines or delay urgent care based on an isolated result.


\section*{Regulatory Status and Authorization Date}
This chapter describes a test category, not a specific commercial product. FDA, CDSCO, EU IVDR, and MHRA approval or registration dates therefore cannot be assigned without the assay manufacturer, product name, instrument, and intended use. For laboratory-developed tests, an individual agency approval date may not exist. See the Preface for the interpretation of regulatory status.

\clearpage
